\section{A stable standard factor and a permuted standard pair}
\label{sec:permuted-standard-pair-cube}

\paragraph{Accepted status and shared provenance.}
D-TD53, reserved question Q2fi, is a joint root/G34 derivation,
written by \texttt{/root/g34\_eightfolds}.  Root and the author
independently reached the same cubic mixture in contemporaneous
messages while studying the surviving two-component standard layer.
The existing stable standard factor supplies the third factor in
arithmetic--geometric mean.  The author sharpens the uniform error
to \((m-1)^{-3/2}\) and supplies the complete proof and constants
below.  The actual induced-cover argument is the accepted D-TD49
interface.  The complete author whole-proof reading passes.
Finite supplementary checking and source/status records are recorded
separately; complete independent review and root acceptance pass.
No new external classification theorem is introduced.

\begin{theorem}[D-TD53]
Let \(C\to\mathbf P^1\) be an actual connected finite
\(J\)-Galois cover with a faithful finite complex coefficient
representation \(V\).  Suppose there are actual homomorphisms
\[
 \begin{gathered}
 x:J\twoheadrightarrow P_n,\qquad P_n\in\{A_n,S_n\},\qquad
 w:J\longrightarrow S_m\wr S_2,\\
 w(J)\supseteq A_m^2,\qquad
 \pi(w(J))=S_2,\qquad n\ge m\ge2^{30}.
 \end{gathered}
 \tag{PSC.1}
\]
Let \(\mathcal W_m\) be the ordinary representation of
\(S_m\wr S_2\) on
\(\operatorname{Std}_m\otimes\operatorname{Std}_m\), with
\(S_2\) exchanging the two tensor factors.  Assume the actual
factorization
\[
 V=(\operatorname{Std}_n\circ x)\otimes
                      (\mathcal W_m\circ w)\otimes R,
 \quad a=m-1,\quad b=n-1,\quad q=\dim R\ge1,
 \quad e=ba^2q,
 \tag{PSC.2}
\]
where \(R\) is the entire remaining representation of J.
If the actual original inertia operators satisfy
\[
                   \sum_i\operatorname{rank}(1-T_i)=2e+6,
 \tag{PSC.3}
\]
then this configuration does not exist.

No maximality, coefficient-field restriction, period-rank
hypothesis, lower bound on the number of original marks, or
support floor on R is required.  No independence of the two
actual quotient maps is assumed.  Irreducibility and primitivity
of V are not needed beyond what a given application may use to
supply (PSC.1)--(PSC.2).
\end{theorem}

\begin{proof}
\subsection*{Actual operators and two genuine cube actions}

The standard character at \(y\in S_N\), with f fixed
letters, is \(f-1\).  Define
\[
 \rho_N(y)=\frac{|f(y)-1|}{N-1}.
 \qquad
 \rho_N(y)^3\le\left(\frac{f(y)}N\right)^3
                  +\frac{\mathbf1_{f(y)=0}}{(N-1)^3}.
 \tag{PSC.4}
\]
Indeed, if \(f\ge1\), then
\((f-1)/(N-1)\le f/N\); the remaining case is exact.
Every finite representation can be made unitary, so normalized
absolute trace is at most one, including on the whole R.

The actual quotient x acts on \(\Omega_n=[n]^3\).
The actual wreath quotient w acts on
\(\Omega_m=\{1,2\}\times[m]^3\), by moving the component
and applying its natural permutation to all three letters.
Write their fixed-point counts as \(F_n(g),F_m(g)\).
Set
\[
 \Phi(g)=\frac{F_n(g)}{3n^3}+\frac{F_m(g)}{3m^3}
       =\frac13\frac{F_n(g)}{n^3}
          +\frac23\frac{F_m(g)}{2m^3},
       \qquad \epsilon=a^{-3/2}.
 \tag{PSC.5}
\]
This is a convex combination of the normalized characters of
two actual permutation actions.  It need not be the normalized
character of one transitive action.

If a standard extraction layer was originally specified up to
scalars, fix its ordinary standard models once.  Perfectness
of the represented \(A_N\) kills any scalar discrepancy on
its core; Schur's lemma makes each compatible local operator
a scalar times the corresponding ordinary \(S_N\) matrix.
The fixed tensor identifications give exactly the ordinary
wreath model \(\mathcal W_m\).  Absorb the remaining scalar
in R.  Since the two fixed standard models and V are genuine
representations of J, the resulting R is a genuine
representation of the same finite J.  This is the accepted
ordinary normalizer interface, not an independent lift of
abstract components.  Odd local permutations and all scalar
phases are included.

\subsection*{A full-power cubic mixture}

Suppose first that g fixes the two components, with local
natural permutations \(y_1,y_2\).  Tensor trace factorization
and arithmetic--geometric mean give
\[
 \frac{|\operatorname{tr}V(g)|}{e}
 \le\rho_n(x(g))\rho_m(y_1)\rho_m(y_2)
 \le\frac{\rho_n(x(g))^3+\rho_m(y_1)^3+\rho_m(y_2)^3}{3}
 \le\Phi(g)+a^{-3}.
 \tag{PSC.6}
\]
The last inequality uses \(b\ge a\).  If all three fixed-letter
counts are positive, the final error is zero.  The mean
inequality is valid although the three standard degrees differ.

If g exchanges components, the elementary tensor-swap trace
identity gives
\(\operatorname{tr}\mathcal W_m(w(g))
 =\operatorname{tr}(\operatorname{Std}_m(y_1y_2))\).
Consequently
\[
                   \frac{|\operatorname{tr}V(g)|}{e}
                     \le\frac{\rho_n(x(g))}{a}.
 \tag{PSC.7}
\]
Here \(F_m(g)=0\).  If \(f(x(g))\ge1\), put
\(u=f(x(g))/n\).  The exact nonnegative identity
\[
 \frac{u^3}{3}-\frac{u}{a}+\frac{2}{3a^{3/2}}
     =\frac{(\sqrt a\,u-1)^2(\sqrt a\,u+2)}{3a^{3/2}}
     \ge0
 \tag{PSC.8}
\]
shows that (PSC.7) is at most \(\Phi(g)+\epsilon\).
If \(f(x(g))=0\), its bound is \(1/(ab)\le a^{-2}\le\epsilon\),
and \(\Phi(g)=0\).  Since \(a^{-3}\le\epsilon\), all cases give
\[
                       \frac{|\operatorname{tr}V(g)|}{e}
                               \le\Phi(g)+\epsilon.
 \tag{PSC.9}
\]
This applies to every element and hence every power of every
original inertia.  In particular, no local cycle-product order
is substituted for the original coefficient order.

\subsection*{At most ten error-bearing original entries}

Call an original entry bad if it exchanges components, or if
it fixes both components and at least one of its three natural
permutations fixes fewer than half of its letters.  A nonbad
entry has only nonbad powers: every originally fixed letter
remains fixed.  Thus every power of a nonbad entry has zero
error in (PSC.9).

For a permutation y on N letters, every eigenvalue in its
permutation module has multiplicity at most the number c of
cycles.  Its standard module therefore has projective support
at least \(N-1-c=\operatorname{ind}(y)-1\).
If \(f(y)<N/2\), the natural index is greater than \(N/4\),
so
\[
 \frac{\nu(\operatorname{Std}_N(y))}{N-1}
       >\frac{N/4-1}{N-1}>\frac15
                              \qquad(N\ge17).
 \tag{PSC.10}
\]
Here \(\nu(A)\) is the dimension minus the largest eigenspace,
and is unchanged by scalar multiplication.  Tensor spectral
convolution cannot decrease this normalized lower bound on
actual moving rank, regardless of the other factors.

For an exchanging entry, tensor conjugation puts its pair
operator in the form \((D\otimes1)S\), with D diagonal and
S the swap.  There are \(a\) constant words, and every unordered
pair of distinct basis letters spans a two-cycle.  Each target
eigenvalue occurs at most once on each such two-cycle.  Its
largest eigenspace therefore has dimension at most
\(a+\binom a2=a(a+1)/2\).  Thus
\[
  \frac{\nu(\mathcal W_m(w(g)))}{a^2}
                 \ge\frac{a-1}{2a}>\frac15.
 \tag{PSC.11}
\]
The diagonal reduction is legitimate because the finite-order
cycle product is diagonalizable; scalar normalizations do not
alter eigenspace dimensions.  The other factors again preserve
the normalized bound.

Every bad original entry therefore costs more than e/5.
As \(e\ge a^3>30\), (PSC.3) allows at most ten of them.
All remaining original entries are retained and incur zero
comparison error, whether or not their projected inertia is
nontrivial.

\subsection*{The original-order Fourier inequality}

Let \(L_i\) be the order of the original coefficient operator
\(T_i\).  Faithfulness makes the orders of both actual
permutation images divide \(L_i\).  Fourier inversion for
its one-eigenspace, (PSC.9), and the zero-error observation give
\[
 \frac{\operatorname{rank}(1-T_i)}e
 \ge\frac13\frac{\operatorname{ind}_{\Omega_n}(g_i)}{n^3}
     +\frac23\frac{\operatorname{ind}_{\Omega_m}(g_i)}{2m^3}
                       -\epsilon\,\mathbf1_{\{i\text{ bad}\}}.
 \tag{PSC.12}
\]
Indeed the average of each fixed-point count over all \(L_i\)
powers is its number of permutation cycles.  The original
scalar phases remain in \(T_i\) and R throughout this argument.

The same proof can be applied after passage to the faithful
coefficient image if that passage is needed in an application.
An element acting trivially on V acts trivially on every
embedded factor algebra; its natural standard and wreath
images are therefore trivial, by projective faithfulness of
the ordinary standard modules.  The two quotient actions
descend to \(C/\ker V\) over the same sphere.  Branch entries
with trivial coefficient image contribute zero moving rank.
One then uses the orders in that faithful image in (PSC.12).

\subsection*{Both cube actions and every possible component-base genus}

Each cube action has exactly five orbits: the diagonal, the
three choices of which two coordinates coincide, and ordered
distinct triples.  They are actual connected quotient covers.
For \(\Omega_n\) this follows from three-transitivity of
\(A_n\).  For \(\Omega_m\), the full \(A_m^2\) gives the
corresponding local transitivity and the component exchange
joins the two blocks.  No product transitivity between x and w
is required.

The accepted numerical sphere theorem HD-GC06 states that an
actual \(A_N\)- or \(S_N\)-Galois sphere cover, \(N\ge10^6\),
has unordered-triple quotient genus greater than \(N^2/288\).
The ordered-triple quotient is a connected degree-six cover
of it, so its genus is greater than \(N^2/48-5\).
Applying actual Riemann--Hurwitz on the five orbits of
\(\Omega_n\) gives
\[
 \frac{\sum_i\operatorname{ind}_{\Omega_n}(g_i)}{n^3}
                 >2+\frac1{24n}-\frac{20}{n^3}.
 \tag{PSC.13}
\]
The repeated-coordinate orbit genera are nonnegative; none is
assigned a cohomology budget.

For completeness, the induced-action argument of D-TD49 must
allow every genus of the component base.  Let H be the actual
index-two component stabilizer, \(X=C/H\), and \(h=g(X)\).
The first component gives an actual quotient of H onto
\(P_m\in\{A_m,S_m\}\).  Its ordered-triple quotient Y has
degree \((m)_3=m(m-1)(m-2)\) over X and is exactly the
distinct-triple orbit cover of \(\Omega_m\).
If \(h=0\), HD-GC06 and the degree-six cover give
\(g(Y)>m^2/48-5\).
If \(h\ge2\), Riemann--Hurwitz gives
\(g(Y)\ge1+(m)_3\), which is more than that bound.
If \(h=1\), the actual \(P_m\)-Galois cover of X cannot
be unramified, since a finite quotient of the topological
fundamental group of a complex elliptic curve is abelian.
At a genuine branch value take a nonidentity natural y,
which fixes at most \(m-2\) letters.  Its index on ordered
distinct triples is at least
\[
 \frac{(m)_3-(m-2)(m-3)(m-4)}2=3(m-2)^2,
 \quad
 g(Y)\ge1+\frac32(m-2)^2>m^2/48-5.
 \tag{PSC.14}
\]
Thus the five connected orbit covers give, in every case,
\[
 \frac{\sum_i\operatorname{ind}_{\Omega_m}(g_i)}{2m^3}
                       >2+\frac1{48m}-\frac{10}{m^3}.
 \tag{PSC.15}
\]
The component stabilizer is a subgroup of the original source
and its quotient cover is actual.  No source for a discarded
factor, projected twelve-mark condition, or genus-zero
assumption on X is used.

\subsection*{The explicit contradiction}

Summing (PSC.12), using at most ten bad entries and
(PSC.13)--(PSC.15), gives
\[
 \begin{split}
 \frac{\sum_i\operatorname{rank}(1-T_i)}e
 &>2+\frac1{72n}+\frac1{72m}
           -\frac{20}{3n^3}-\frac{20}{3m^3}-10a^{-3/2}\\
 &\ge2+\frac1{72m}-\frac{40}{3m^3}-10a^{-3/2}.
 \end{split}
 \tag{PSC.16}
\]
Since \(m\ge2^{30}\), we have \(a\ge m/2\) and
\(\sqrt a\ge2^{14}\), whence
\(10a^{-3/2}\le5/(4096m)\).
The exact comparison \(1/72-5/4096>1/80\) now shows that
(PSC.16) is greater than
\(2+1/(80m)-40/(3m^3)\).
On the other hand \(e\ge a^3\ge m^3/8\), so \(6/e\le48/m^3\).
Finally \(3m^2>14720\) implies
\[
                   \frac1{80m}-\frac{40}{3m^3}
                        >\frac{48}{m^3}\ge\frac6e.
 \tag{PSC.17}
\]
This contradicts (PSC.3) and proves D-TD53.
\end{proof}

\paragraph{Application boundary.}
In the maximal-standard setting, a separately established
complete decomposition supplies the stable \(\operatorname{Std}_n\)
and a two-component ordinary standard layer with actual
\(A_m^2\) and component exchange.  The original whole
complement supplies R.  If \(m\ge2^{30}\), D-TD53 excludes
that configuration without the near-maximum ratio or its
whole-partner floor.  This result does not by itself prove
that a different extraction layer is standard, supply the
two quotient maps, or impose a new source budget on R.
It is separate from the same-degree three-or-more-component
argument and the two individually stable standard theorem.
